% ============================================================
% sindicancia-procedimento.tex
% Estado processual e validacoes da EB10-IG-09.001 (2a ed., 2024)
% Carregado por sindicancia.cls com \makeatletter ativo.
% ============================================================

% ------------------------------------------------------------
% 1. REGIME NORMATIVO, RITO, NATUREZA E PARTE PRINCIPAL
% ------------------------------------------------------------
\gdef\@sindDataPublicacaoPortariaISO{}
\gdef\@sindDataRecebimentoPortariaISO{}
\ifsind@sumario
  \gdef\@sindRito{sumario}
\else
  \gdef\@sindRito{ordinario}
\fi
\gdef\@sindNatureza{investigatoria}
\gdef\@sindPartePrincipal{nenhuma}
\gdef\@sindClassificacao{ostensiva}
\gdef\@sindSigiloAutoridade{}
\gdef\@sindSigiloFundamento{}
\gdef\@sindSigiloData{}
\gdef\@sindHistoricoTransicoes{}

\newcommand{\dataPublicacaoPortariaISO}[1]{\gdef\@sindDataPublicacaoPortariaISO{#1}}
\newcommand{\dataRecebimentoPortariaISO}[1]{\gdef\@sindDataRecebimentoPortariaISO{#1}}

\newcommand{\ritoSindicancia}[1]{%
  \ifstrequal{#1}{ordinario}{\gdef\@sindRito{ordinario}}{%
    \ifstrequal{#1}{sumario}{\gdef\@sindRito{sumario}}{%
      \ClassError{sindicancia}{Rito invalido: #1}{Use ordinario ou sumario.}%
    }%
  }%
}

\newcommand{\naturezaSindicancia}[1]{%
  \ifstrequal{#1}{investigatoria}{\gdef\@sindNatureza{investigatoria}}{%
    \ifstrequal{#1}{processual}{\gdef\@sindNatureza{processual}}{%
      \ClassError{sindicancia}{Natureza invalida: #1}{Use investigatoria ou processual.}%
    }%
  }%
}

\newcommand{\partePrincipalSindicancia}[1]{%
  \ifstrequal{#1}{nenhuma}{\gdef\@sindPartePrincipal{nenhuma}}{%
    \ifstrequal{#1}{interessado}{\gdef\@sindPartePrincipal{interessado}}{%
      \ifstrequal{#1}{sindicado}{\gdef\@sindPartePrincipal{sindicado}}{%
        \ClassError{sindicancia}{Parte principal invalida: #1}{Use nenhuma, interessado ou sindicado.}%
      }%
    }%
  }%
}

\newcommand{\classificacaoSindicancia}[1]{%
  \ifstrequal{#1}{ostensiva}{\gdef\@sindClassificacao{ostensiva}}{%
    \ifstrequal{#1}{sigilosa}{\gdef\@sindClassificacao{sigilosa}}{%
      \ClassError{sindicancia}{Classificacao invalida: #1}{Use ostensiva ou sigilosa.}%
    }%
  }%
}
\newcommand{\sigiloAutoridade}[1]{\gdef\@sindSigiloAutoridade{#1}}
\newcommand{\sigiloFundamento}[1]{\gdef\@sindSigiloFundamento{#1}}
\newcommand{\sigiloData}[1]{\gdef\@sindSigiloData{#1}}

\newcommand{\getRitoSindicancia}{\@sindRito}
\newcommand{\getNaturezaSindicancia}{\@sindNatureza}
\newcommand{\getPartePrincipalSindicancia}{\@sindPartePrincipal}
\newcommand{\getClassificacaoSindicancia}{\@sindClassificacao}
\newcommand{\getSigiloAutoridade}{\@sindSigiloAutoridade}
\newcommand{\getSigiloFundamento}{\@sindSigiloFundamento}
\newcommand{\getSigiloData}{\@sindSigiloData}
\newcommand{\getHistoricoTransicoes}{\@sindHistoricoTransicoes}
\renewcommand{\getTipoSindicancia}{\@sindRito}

\newcommand{\registrarTransicao}[5]{%
  \g@addto@macro\@sindHistoricoTransicoes{%
    \item \textbf{#1}: #2, de #3 para #4. Fundamento: #5.%
  }%
}

\newcommand{\converterParaProcessual}[2]{%
  \ifdefstring{\@sindNatureza}{investigatoria}{}{%
    \ClassError{sindicancia}{Conversao processual invalida}{A conversao exige natureza investigatoria anterior.}%
  }%
  \ValidarValorObrigatorio{#1}{data da conversao para natureza processual}%
  \ValidarValorObrigatorio{#2}{fundamento da conversao para natureza processual}%
  \registrarTransicao{#1}{natureza}{\@sindNatureza}{processual}{#2}%
  \gdef\@sindNatureza{processual}%
  \gdef\@sindPartePrincipal{sindicado}%
}

\newcommand{\converterRitoParaOrdinario}[2]{%
  \ifdefstring{\@sindRito}{sumario}{}{%
    \ClassError{sindicancia}{Mudanca de rito invalida}{A mudanca para ordinario exige rito sumario anterior.}%
  }%
  \ValidarValorObrigatorio{#1}{data da conversao para rito ordinario}%
  \ValidarValorObrigatorio{#2}{fundamento da conversao para rito ordinario}%
  \registrarTransicao{#1}{rito}{\@sindRito}{ordinario}{#2}%
  \gdef\@sindRito{ordinario}%
}

\newcommand{\ValidarRegimeNormativo}{%
  \ValidarValorObrigatorio{\@sindDataPublicacaoPortariaISO}{data de publicacao da portaria no formato AAAAMMDD}%
  \ifnum\@sindDataPublicacaoPortariaISO<20250101\relax
    \ClassError{sindicancia}
      {Portaria anterior a 1 de janeiro de 2025}
      {O caso e regido pela Portaria C Ex 107/2012 e nao por esta versao do projeto.}%
  \fi
  \ValidarValorObrigatorio{\@sindDataRecebimentoPortariaISO}{data de recebimento da portaria no formato AAAAMMDD}%
}

\newcommand{\ValidarEstadoSindicancia}{%
  \ifdefstring{\@sindNatureza}{processual}{%
    \ifdefstring{\@sindPartePrincipal}{sindicado}{}{%
      \ClassError{sindicancia}{Estado processual sem sindicado}{Defina a parte principal como sindicado.}%
    }%
    \CampoObrigatorio{sindicadoNome}{nome do sindicado}%
    \CampoObrigatorio{sindicadoPosto}{posto, graduacao ou qualificacao do sindicado}%
  }{}%
  \ifdefstring{\@sindPartePrincipal}{interessado}{%
    \CampoObrigatorio{interessadoNome}{nome do interessado}%
    \CampoObrigatorio{interessadoPosto}{posto, graduacao ou qualificacao do interessado}%
  }{}%
  \ifdefstring{\@sindClassificacao}{sigilosa}{%
    \ValidarValorObrigatorio{\@sindSigiloAutoridade}{autoridade que classificou a sindicancia}%
    \ValidarValorObrigatorio{\@sindSigiloFundamento}{fundamento da classificacao sigilosa}%
    \ValidarValorObrigatorio{\@sindSigiloData}{data da classificacao sigilosa}%
  }{}%
}

\newcommand{\sind@cabecalhoClassificacao}{%
  \ifdefstring{\@sindClassificacao}{sigilosa}{%
    \vspace{0.3em}\textbf{\MakeUppercase{Documento sigiloso -- acesso controlado}}\\%
  }{}%
}

% ------------------------------------------------------------
% 2. DATAS ESPECIFICAS E ORDEM CRONOLOGICA
% ------------------------------------------------------------
\gdef\@documentoChaveCronologica{}
\gdef\@ultimaChaveCronologica{0}
\newcommand{\chaveCronologica}[1]{\gdef\@documentoChaveCronologica{#1}}
\newcommand{\getChaveCronologica}{\@documentoChaveCronologica}

\newcommand{\ValidarDataDocumentoEspecifica}[1]{%
  \ValidarValorObrigatorio{\@documentoDia}{dia de #1}%
  \ValidarValorObrigatorio{\@documentoMes}{mes de #1}%
  \ValidarValorObrigatorio{\@documentoAno}{ano de #1}%
  \ValidarValorObrigatorio{\@documentoCidade}{cidade de #1}%
  \ValidarValorObrigatorio{\@documentoChaveCronologica}{chave cronologica AAAAMMDD de #1}%
}

\newcommand{\ValidarCronologiaDocumento}[1]{%
  \ifnum\@documentoChaveCronologica<\@ultimaChaveCronologica\relax
    \ClassError{sindicancia}
      {Documento fora da ordem cronologica: #1}
      {Reordene os documentos no main.tex ou corrija a chave cronologica AAAAMMDD.}%
  \fi
  \xdef\@ultimaChaveCronologica{\@documentoChaveCronologica}%
}

\newcommand{\ValidarDataDocumentoParaModelo}[1]{%
  \ifstrequal{#1}{capa}{}{%
    \ValidarDataDocumentoEspecifica{#1}%
    \ValidarCronologiaDocumento{#1}%
  }%
}

\apptocmd{\limparDataDocumento}{\gdef\@documentoChaveCronologica{}}{}{}

% ------------------------------------------------------------
% 3. FLEXIBILIDADE DE CONTEUDO
% ------------------------------------------------------------
\gdef\@modoDocumento{padrao}
\gdef\@textoDocumentoPersonalizado{}
\gdef\@textoComplementarDocumento{}

\newcommand{\modoDocumento}[1]{\gdef\@modoDocumento{#1}}
\newcommand{\textoDocumentoPersonalizado}[1]{\gdef\@textoDocumentoPersonalizado{#1}}
\newcommand{\textoComplementarDocumento}[1]{\gdef\@textoComplementarDocumento{#1}}

\newcommand{\sind@limparFlexibilidade}{%
  \gdef\@modoDocumento{padrao}%
  \gdef\@textoDocumentoPersonalizado{}%
  \gdef\@textoComplementarDocumento{}%
}

\newcommand{\ValidarModoDocumento}{%
  \ifdefstring{\@modoDocumento}{padrao}{}{%
    \ifdefstring{\@modoDocumento}{adaptado}{}{%
      \ifdefstring{\@modoDocumento}{personalizado}{%
        \ValidarValorObrigatorio{\@textoDocumentoPersonalizado}{texto personalizado do documento}%
      }{%
        \ClassError{sindicancia}{Modo de documento invalido: \@modoDocumento}{Use padrao, adaptado ou personalizado.}%
      }%
    }%
  }%
}

\newcommand{\ConteudoDocumentoFlexivel}[1]{%
  \ValidarModoDocumento
  \ifdefstring{\@modoDocumento}{personalizado}{%
    \@textoDocumentoPersonalizado
  }{%
    #1%
    \ifdefstring{\@modoDocumento}{adaptado}{%
      \SeCampoPreenchido{textoComplementarDocumento}{\par\vspace{0.8em}\@textoComplementarDocumento}%
    }{}%
  }%
}

\apptocmd{\limparDataDocumento}{\sind@limparFlexibilidade}{}{}

% ------------------------------------------------------------
% 4. CADASTRO ESTRUTURADO DE PARTICIPANTES
% ------------------------------------------------------------
\gdef\@participantesIds{}
\newcommand{\registrarParticipante}[6]{%
  \ifcsdef{sind@participante@#1@nome}{%
    \ClassError{sindicancia}{Participante duplicado: #1}{Use um identificador unico para cada participante.}%
  }{}%
  \expandafter\gdef\csname sind@participante@#1@papel\endcsname{#2}%
  \expandafter\gdef\csname sind@participante@#1@nome\endcsname{#3}%
  \expandafter\gdef\csname sind@participante@#1@qualificacao\endcsname{#4}%
  \expandafter\gdef\csname sind@participante@#1@vinculo\endcsname{#5}%
  \expandafter\gdef\csname sind@participante@#1@observacoes\endcsname{#6}%
  \listgadd{\@participantesIds}{#1}%
}
\newcommand{\getParticipantePapel}[1]{\csname sind@participante@#1@papel\endcsname}
\newcommand{\getParticipanteNome}[1]{\csname sind@participante@#1@nome\endcsname}
\newcommand{\getParticipanteQualificacao}[1]{\csname sind@participante@#1@qualificacao\endcsname}
\newcommand{\getParticipanteVinculo}[1]{\csname sind@participante@#1@vinculo\endcsname}
\newcommand{\getParticipantesIds}{\@participantesIds}

\gdef\@quantTestemunhasDenunciante{0}
\gdef\@quantTestemunhasSindicado{0}
\newcommand{\quantidadeTestemunhasDenunciante}[1]{\gdef\@quantTestemunhasDenunciante{#1}}
\newcommand{\quantidadeTestemunhasSindicado}[1]{\gdef\@quantTestemunhasSindicado{#1}}
\newcommand{\ValidarLimiteTestemunhas}{%
  \ifnum\@quantTestemunhasDenunciante>3\relax
    \ClassError{sindicancia}{Mais de tres testemunhas indicadas pelo denunciante ou ofendido}{O art. 17 limita essa indicacao a tres testemunhas.}%
  \fi
  \ifnum\@quantTestemunhasSindicado>3\relax
    \ClassError{sindicancia}{Mais de tres testemunhas indicadas pelo sindicado}{O art. 17 limita essa indicacao a tres testemunhas.}%
  \fi
}

% ------------------------------------------------------------
% 5. REGRAS DE INQUIRICAO
% ------------------------------------------------------------
\gdef\@testemunhaLocalTrabalho{}
\gdef\@testemunhaParentesco{}
\gdef\@testemunhaDispensaCompromisso{}
\gdef\@testemunhaSegredoProfissional{}
\gdef\@testemunhaDesobrigadaPor{}
\gdef\@modalidadeInquiricao{presencial}
\gdef\@regulamentoAtoRemoto{}
\gdef\@inquiricaoInicioMinutos{}
\gdef\@inquiricaoFimMinutos{}
\gdef\@inquiricaoPausaMinutos{0}
\gdef\@inquiricaoJustificativaHorario{}
\gdef\@inquiricaoContinuacao{}
\gdef\@assinaturaRogoNome{}
\gdef\@assinaturaRogoMotivo{}
\gdef\@assinaturaRogoTestemunhaA{}
\gdef\@assinaturaRogoTestemunhaB{}
\gdef\@sindicadoMenor{nao}
\gdef\@responsavelMenorNome{}
\gdef\@responsavelMenorQualificacao{}
\gdef\@manifestacaoDefesaInquiricao{}

\newcommand{\sind@limparRegrasInquiricao}{%
  \gdef\@testemunhaLocalTrabalho{}%
  \gdef\@testemunhaParentesco{}%
  \gdef\@testemunhaDispensaCompromisso{}%
  \gdef\@testemunhaSegredoProfissional{}%
  \gdef\@testemunhaDesobrigadaPor{}%
  \gdef\@modalidadeInquiricao{presencial}%
  \gdef\@regulamentoAtoRemoto{}%
  \gdef\@inquiricaoInicioMinutos{}%
  \gdef\@inquiricaoFimMinutos{}%
  \gdef\@inquiricaoPausaMinutos{0}%
  \gdef\@inquiricaoJustificativaHorario{}%
  \gdef\@inquiricaoContinuacao{}%
  \gdef\@assinaturaRogoNome{}%
  \gdef\@assinaturaRogoMotivo{}%
  \gdef\@assinaturaRogoTestemunhaA{}%
  \gdef\@assinaturaRogoTestemunhaB{}%
  \gdef\@sindicadoMenor{nao}%
  \gdef\@responsavelMenorNome{}%
  \gdef\@responsavelMenorQualificacao{}%
  \gdef\@manifestacaoDefesaInquiricao{}%
}
\apptocmd{\novaInquiricaoTestemunha}{\sind@limparRegrasInquiricao}{}{}
\apptocmd{\novaInquiricaoSindicado}{\sind@limparRegrasInquiricao}{}{}

\newcommand{\testemunhaLocalTrabalho}[1]{\gdef\@testemunhaLocalTrabalho{#1}}
\newcommand{\testemunhaParentesco}[1]{\gdef\@testemunhaParentesco{#1}}
\newcommand{\testemunhaDispensaCompromisso}[1]{\gdef\@testemunhaDispensaCompromisso{#1}}
\newcommand{\testemunhaSegredoProfissional}[1]{\gdef\@testemunhaSegredoProfissional{#1}}
\newcommand{\testemunhaDesobrigadaPor}[1]{\gdef\@testemunhaDesobrigadaPor{#1}}
\newcommand{\modalidadeInquiricao}[1]{\gdef\@modalidadeInquiricao{#1}}
\newcommand{\regulamentoAtoRemoto}[1]{\gdef\@regulamentoAtoRemoto{#1}}
\newcommand{\inquiricaoHorarioMinutos}[2]{\gdef\@inquiricaoInicioMinutos{#1}\gdef\@inquiricaoFimMinutos{#2}}
\newcommand{\inquiricaoPausaMinutos}[1]{\gdef\@inquiricaoPausaMinutos{#1}}
\newcommand{\inquiricaoJustificativaHorario}[1]{\gdef\@inquiricaoJustificativaHorario{#1}}
\newcommand{\inquiricaoContinuacao}[1]{\gdef\@inquiricaoContinuacao{#1}}
\newcommand{\assinaturaRogo}[4]{%
  \gdef\@assinaturaRogoNome{#1}%
  \gdef\@assinaturaRogoMotivo{#2}%
  \gdef\@assinaturaRogoTestemunhaA{#3}%
  \gdef\@assinaturaRogoTestemunhaB{#4}%
}
\newcommand{\sindicadoMenor}[1]{\gdef\@sindicadoMenor{#1}}
\newcommand{\responsavelMenor}[2]{\gdef\@responsavelMenorNome{#1}\gdef\@responsavelMenorQualificacao{#2}}
\newcommand{\manifestacaoDefesaInquiricao}[1]{\gdef\@manifestacaoDefesaInquiricao{#1}}
\newcommand{\getTestemunhaLocalTrabalho}{\@testemunhaLocalTrabalho}
\newcommand{\getTestemunhaParentesco}{\@testemunhaParentesco}
\newcommand{\getTestemunhaDispensaCompromisso}{\@testemunhaDispensaCompromisso}
\newcommand{\getModalidadeInquiricao}{\@modalidadeInquiricao}
\newcommand{\getRegulamentoAtoRemoto}{\@regulamentoAtoRemoto}
\newcommand{\getManifestacaoDefesaInquiricao}{\@manifestacaoDefesaInquiricao}

\newcommand{\TextoModalidadeInquiricao}{%
  \ifdefstring{\@modalidadeInquiricao}{presencial}{presencialmente}{%
    \ifdefstring{\@modalidadeInquiricao}{videoconferencia}{por videoconferência}{%
      \ifdefstring{\@modalidadeInquiricao}{telepresencial}{de forma telepresencial}{%
        por carta precatória%
      }%
    }%
  }%
}

\newcommand{\TextoCompromissoTestemunha}{%
  \ifcsempty{@testemunhaDispensaCompromisso}{%
    prestou o compromisso de dizer a verdade sobre o que souber e lhe for perguntado%
  }{%
    não prestou compromisso, em razão de \@testemunhaDispensaCompromisso%
  }%
}
\newcommand{\TextoSegredoProfissionalTestemunha}{%
  \ifcsempty{@testemunhaSegredoProfissional}{}{%
    Declarou estar sujeita a segredo profissional quanto a \@testemunhaSegredoProfissional\ e foi desobrigada por \@testemunhaDesobrigadaPor. %
  }%
}

\newcommand{\ValidarRegrasInquiricao}{%
  \ifdefstring{\@modalidadeInquiricao}{presencial}{}{%
    \ifdefstring{\@modalidadeInquiricao}{videoconferencia}{}{%
      \ifdefstring{\@modalidadeInquiricao}{telepresencial}{}{%
        \ifdefstring{\@modalidadeInquiricao}{carta-precatoria}{}{%
          \ClassError{sindicancia}{Modalidade de inquiricao invalida}{Use presencial, videoconferencia, telepresencial ou carta-precatoria.}%
        }%
      }%
    }%
  }%
  \ifdefstring{\@modalidadeInquiricao}{videoconferencia}{\ValidarValorObrigatorio{\@regulamentoAtoRemoto}{regulamento do ato remoto}}{}%
  \ifdefstring{\@modalidadeInquiricao}{telepresencial}{\ValidarValorObrigatorio{\@regulamentoAtoRemoto}{regulamento do ato remoto}}{}%
  \ValidarValorObrigatorio{\@inquiricaoInicioMinutos}{horario inicial da inquiricao em minutos}%
  \ValidarValorObrigatorio{\@inquiricaoFimMinutos}{horario final da inquiricao em minutos}%
  \ifnum\@inquiricaoFimMinutos<\@inquiricaoInicioMinutos\relax
    \ClassError{sindicancia}{Horario final anterior ao inicial}{Corrija os horarios da inquiricao.}%
  \fi
  \ifnum\@inquiricaoInicioMinutos<480\relax
    \ValidarValorObrigatorio{\@inquiricaoJustificativaHorario}{justificativa para ato antes das 8h}%
  \fi
  \ifnum\@inquiricaoFimMinutos>1080\relax
    \ValidarValorObrigatorio{\@inquiricaoJustificativaHorario}{justificativa para ato depois das 18h}%
  \fi
  \@tempcnta=\@inquiricaoFimMinutos\relax
  \advance\@tempcnta by -\@inquiricaoInicioMinutos\relax
  \ifnum\@tempcnta>240\relax
    \ifnum\@inquiricaoPausaMinutos<30\relax
      \ifcsempty{@inquiricaoContinuacao}{%
        \ClassError{sindicancia}{Inquiricao superior a quatro horas sem pausa ou continuacao}{Registre pausa minima de 30 minutos ou a continuacao em outro dia.}%
      }{}%
    \fi
  \fi
  \ifcsempty{@assinaturaRogoNome}{}{%
    \ValidarValorObrigatorio{\@assinaturaRogoMotivo}{motivo da assinatura a rogo}%
    \ValidarValorObrigatorio{\@assinaturaRogoTestemunhaA}{primeira testemunha da assinatura a rogo}%
    \ValidarValorObrigatorio{\@assinaturaRogoTestemunhaB}{segunda testemunha da assinatura a rogo}%
  }%
  \ifcsempty{@testemunhaInquiricaoBNome}{}{%
    \CampoObrigatorio{testemunhaInquiricaoNome}{primeira testemunha instrumentaria}%
  }%
  \ifdefstring{\@sindicadoMenor}{sim}{%
    \ValidarValorObrigatorio{\@responsavelMenorNome}{responsavel pelo sindicado menor}%
    \ValidarValorObrigatorio{\@responsavelMenorQualificacao}{qualificacao do responsavel pelo sindicado menor}%
  }{}%
  \ifcsempty{@testemunhaSegredoProfissional}{}{%
    \ValidarValorObrigatorio{\@testemunhaDesobrigadaPor}{parte que desobrigou a testemunha do segredo profissional}%
  }%
}

\newcommand{\AssinaturasRogoSeHouver}{%
  \ifcsempty{@assinaturaRogoNome}{}{%
    \assinaturaCentral{\@assinaturaRogoNome}{Assinatura a rogo -- \@assinaturaRogoMotivo}%
    \assinaturaCentral{\@assinaturaRogoTestemunhaA}{Testemunha da assinatura a rogo}%
    \assinaturaCentral{\@assinaturaRogoTestemunhaB}{Testemunha da assinatura a rogo}%
  }%
}

\newcommand{\ResponsavelMenorSeHouver}{%
  \ifdefstring{\@sindicadoMenor}{sim}{%
    \assinaturaCentral{\@responsavelMenorNome}{\@responsavelMenorQualificacao -- responsável legal}%
  }{}%
}
\newcommand{\TextoResponsavelMenorSeHouver}{%
  \ifdefstring{\@sindicadoMenor}{sim}{Compareceu também seu responsável legal. }{}%
}

% ------------------------------------------------------------
% 6. NOTIFICACOES, EXPEDIENTES E CARTA PRECATORIA
% ------------------------------------------------------------
\gdef\@notificacaoCanal{direta}
\gdef\@notificacaoAntecedenciaDiasUteis{3}
\gdef\@notificacaoFatosEspecificos{}
\gdef\@notificacaoFormaCiencia{}
\newcommand{\notificacaoCanal}[1]{\gdef\@notificacaoCanal{#1}}
\newcommand{\notificacaoAntecedenciaDiasUteis}[1]{\gdef\@notificacaoAntecedenciaDiasUteis{#1}}
\newcommand{\notificacaoFatosEspecificos}[1]{\gdef\@notificacaoFatosEspecificos{#1}}
\newcommand{\notificacaoFormaCiencia}[1]{\gdef\@notificacaoFormaCiencia{#1}}
\newcommand{\getNotificacaoCanal}{\@notificacaoCanal}
\newcommand{\getNotificacaoAntecedenciaDiasUteis}{\@notificacaoAntecedenciaDiasUteis}
\newcommand{\getNotificacaoFatosEspecificos}{\@notificacaoFatosEspecificos}
\newcommand{\getNotificacaoFormaCiencia}{\@notificacaoFormaCiencia}
\newcommand{\ValidarRegrasNotificacao}{%
  \ifnum\@notificacaoAntecedenciaDiasUteis<3\relax
    \ClassError{sindicancia}{Antecedencia de notificacao inferior a tres dias uteis}{O art. 32 exige antecedencia minima de tres dias uteis para diligencias.}%
  \fi
  \ValidarValorObrigatorio{\@notificacaoFatosEspecificos}{descricao especifica dos fatos comunicados}%
  \ValidarValorObrigatorio{\@notificacaoFormaCiencia}{forma de ciencia da notificacao}%
}

\gdef\@expedienteFinalidade{}
\gdef\@expedientePrioridade{normal}
\newcommand{\expedienteFinalidade}[1]{\gdef\@expedienteFinalidade{#1}}
\newcommand{\expedientePrioridade}[1]{\gdef\@expedientePrioridade{#1}}
\newcommand{\getExpedienteFinalidade}{\@expedienteFinalidade}
\newcommand{\getExpedientePrioridade}{\@expedientePrioridade}
\newcommand{\ValidarRegrasExpediente}{%
  \ValidarValorObrigatorio{\@expedienteFinalidade}{finalidade do pedido ou requisicao}%
  \ifdefstring{\@expedientePrioridade}{normal}{}{%
    \ifdefstring{\@expedientePrioridade}{urgente}{}{%
      \ifdefstring{\@expedientePrioridade}{urgentissima}{}{%
        \ClassError{sindicancia}{Prioridade de expediente invalida}{Use normal, urgente ou urgentissima.}%
      }%
    }%
  }%
}

\gdef\@cartaPrecatoriaQuesitosDefesa{}
\gdef\@cartaPrecatoriaNotificacaoQuesitos{}
\newcommand{\cartaPrecatoriaQuesitosDefesa}[1]{\gdef\@cartaPrecatoriaQuesitosDefesa{#1}}
\newcommand{\cartaPrecatoriaNotificacaoQuesitos}[1]{\gdef\@cartaPrecatoriaNotificacaoQuesitos{#1}}
\newcommand{\getCartaPrecatoriaQuesitosDefesa}{\@cartaPrecatoriaQuesitosDefesa}
\newcommand{\getCartaPrecatoriaNotificacaoQuesitos}{\@cartaPrecatoriaNotificacaoQuesitos}
\newcommand{\ValidarRegrasCartaPrecatoria}{%
  \ValidarRegrasExpediente
  \ifdefstring{\@sindPartePrincipal}{sindicado}{%
    \ValidarValorObrigatorio{\@cartaPrecatoriaNotificacaoQuesitos}{notificacao do sindicado para apresentacao de quesitos}%
  }{}%
}

% ------------------------------------------------------------
% 7. ACAREACAO COM PARTICIPANTES DINAMICOS
% ------------------------------------------------------------
\newcounter{sindAcareados}
\gdef\@acareacaoDeclaracoesDinamicas{}
\gdef\@acareacaoAssinaturasDinamicas{}
\newcommand{\novaAcareacaoDinamica}{%
  \setcounter{sindAcareados}{0}%
  \gdef\@acareacaoDeclaracoesDinamicas{}%
  \gdef\@acareacaoAssinaturasDinamicas{}%
}
\newcommand{\adicionarParticipanteAcareacao}[3]{%
  \stepcounter{sindAcareados}%
  \g@addto@macro\@acareacaoDeclaracoesDinamicas{\par\textbf{#1 (#2):} #3}%
  \g@addto@macro\@acareacaoAssinaturasDinamicas{\assinaturaCentral{#1}{#2}}%
}
\newcommand{\ValidarAcareacaoDinamica}{%
  \ifnum\value{sindAcareados}<2\relax
    \ClassError{sindicancia}{Acareacao com menos de dois participantes}{Registre ao menos dois depoentes com declaracoes divergentes.}%
  \fi
}
\newcommand{\getAcareacaoDeclaracoesDinamicas}{\@acareacaoDeclaracoesDinamicas}
\newcommand{\getAcareacaoAssinaturasDinamicas}{\@acareacaoAssinaturasDinamicas}

% ------------------------------------------------------------
% 8. CAMPOS DOS ANEXOS A, B, Z E AA
% ------------------------------------------------------------
\newcommand{\AvisoModeloApoioJuridico}[1]{%
  \ClassWarningNoLine{sindicancia}{Anexo #1 mantido apenas como referencia; sua elaboracao compete a Secao de Apoio Juridico, nao ao sindicante}%
}

\gdef\@portariaNumero{}
\gdef\@portariaAnexoOrigem{}
\gdef\@portariaSinteseFatos{}
\gdef\@portariaPrazoDias{30}
\gdef\@portariaEscrivao{}
\gdef\@portariaPublicacaoBI{}
\gdef\@portariaAnonimaDocumentoVerificacao{}
\gdef\@portariaAnonimaEnvolvido{}
\gdef\@portariaAnonimaOM{}
\gdef\@portariaAnonimaIrregularidade{}
\gdef\@portariaAnonimaInteressePublico{}
\newif\ifsind@denunciaanonimanaojuntada
\sind@denunciaanonimanaojuntadafalse

\newcommand{\portariaNumero}[1]{\gdef\@portariaNumero{#1}}
\newcommand{\portariaAnexoOrigem}[1]{\gdef\@portariaAnexoOrigem{#1}}
\newcommand{\portariaSinteseFatos}[1]{\gdef\@portariaSinteseFatos{#1}}
\newcommand{\portariaPrazoDias}[1]{\gdef\@portariaPrazoDias{#1}}
\newcommand{\portariaEscrivao}[1]{\gdef\@portariaEscrivao{#1}}
\newcommand{\portariaPublicacaoBI}[1]{\gdef\@portariaPublicacaoBI{#1}}
\newcommand{\portariaAnonimaDocumentoVerificacao}[1]{\gdef\@portariaAnonimaDocumentoVerificacao{#1}}
\newcommand{\portariaAnonimaEnvolvido}[1]{\gdef\@portariaAnonimaEnvolvido{#1}}
\newcommand{\portariaAnonimaOM}[1]{\gdef\@portariaAnonimaOM{#1}}
\newcommand{\portariaAnonimaIrregularidade}[1]{\gdef\@portariaAnonimaIrregularidade{#1}}
\newcommand{\portariaAnonimaInteressePublico}[1]{\gdef\@portariaAnonimaInteressePublico{#1}}
\newcommand{\confirmarDenunciaAnonimaNaoJuntada}{\sind@denunciaanonimanaojuntadatrue}
\newcommand{\getPortariaNumero}{\@portariaNumero}
\newcommand{\getPortariaAnexoOrigem}{\@portariaAnexoOrigem}
\newcommand{\getPortariaSinteseFatos}{\@portariaSinteseFatos}
\newcommand{\getPortariaPrazoDias}{\@portariaPrazoDias}
\newcommand{\getPortariaEscrivao}{\@portariaEscrivao}
\newcommand{\getPortariaPublicacaoBI}{\@portariaPublicacaoBI}
\newcommand{\getPortariaAnonimaDocumentoVerificacao}{\@portariaAnonimaDocumentoVerificacao}
\newcommand{\getPortariaAnonimaEnvolvido}{\@portariaAnonimaEnvolvido}
\newcommand{\getPortariaAnonimaOM}{\@portariaAnonimaOM}
\newcommand{\getPortariaAnonimaIrregularidade}{\@portariaAnonimaIrregularidade}
\newcommand{\getPortariaAnonimaInteressePublico}{\@portariaAnonimaInteressePublico}

\newcommand{\ValidarPortariaInstauracao}{%
  \CampoObrigatorio{autoridadeNome}{nome da autoridade instauradora}%
  \CampoObrigatorio{autoridadePosto}{posto da autoridade instauradora}%
  \CampoObrigatorio{autoridadeFuncao}{funcao da autoridade instauradora}%
  \CampoObrigatorio{sindicanteNome}{nome do sindicante}%
  \CampoObrigatorio{sindicantePosto}{posto ou graduacao do sindicante}%
  \ValidarValorObrigatorio{\@portariaNumero}{numero da portaria}%
  \ValidarValorObrigatorio{\@portariaAnexoOrigem}{documento que motivou a sindicancia}%
  \ValidarValorObrigatorio{\@portariaSinteseFatos}{sintese dos fatos}%
  \ValidarValorObrigatorio{\@portariaPublicacaoBI}{publicacao da portaria em BI}%
}

\newcommand{\ValidarPortariaDenunciaAnonima}{%
  \ValidarPortariaInstauracao
  \ValidarValorObrigatorio{\@portariaAnonimaDocumentoVerificacao}{documento produzido na verificacao de plausibilidade}%
  \ValidarValorObrigatorio{\@portariaAnonimaIrregularidade}{irregularidade verificada}%
  \ValidarValorObrigatorio{\@portariaAnonimaInteressePublico}{interesse publico na apuracao}%
  \ifsind@denunciaanonimanaojuntada\else
    \ClassError{sindicancia}{Confirmacao ausente sobre denuncia anonima}{Use \string\confirmarDenunciaAnonimaNaoJuntada depois de verificar que a peca anonima nao integra os autos.}%
  \fi
}

\gdef\@solucaoTipo{acolhe}
\gdef\@solucaoParecer{}
\gdef\@solucaoFundamentos{}
\gdef\@solucaoProvidencias{}
\gdef\@solucaoPublicacaoBI{}
\gdef\@solucaoNotificados{}
\newcommand{\solucaoTipo}[1]{\gdef\@solucaoTipo{#1}}
\newcommand{\solucaoParecer}[1]{\gdef\@solucaoParecer{#1}}
\newcommand{\solucaoFundamentos}[1]{\gdef\@solucaoFundamentos{#1}}
\newcommand{\solucaoProvidencias}[1]{\gdef\@solucaoProvidencias{#1}}
\newcommand{\solucaoPublicacaoBI}[1]{\gdef\@solucaoPublicacaoBI{#1}}
\newcommand{\solucaoNotificados}[1]{\gdef\@solucaoNotificados{#1}}
\newcommand{\getSolucaoTipo}{\@solucaoTipo}
\newcommand{\getSolucaoParecer}{\@solucaoParecer}
\newcommand{\getSolucaoFundamentos}{\@solucaoFundamentos}
\newcommand{\getSolucaoProvidencias}{\@solucaoProvidencias}
\newcommand{\getSolucaoPublicacaoBI}{\@solucaoPublicacaoBI}
\newcommand{\getSolucaoNotificados}{\@solucaoNotificados}
\newcommand{\ValidarSolucao}{%
  \CampoObrigatorio{autoridadeNome}{nome da autoridade instauradora}%
  \CampoObrigatorio{autoridadePosto}{posto da autoridade instauradora}%
  \ValidarValorObrigatorio{\@solucaoParecer}{parecer acolhido, parcialmente acolhido ou rejeitado}%
  \ValidarValorObrigatorio{\@solucaoFundamentos}{fundamentos da solucao}%
  \ValidarValorObrigatorio{\@solucaoProvidencias}{providencias da solucao}%
  \ValidarValorObrigatorio{\@solucaoPublicacaoBI}{publicacao da solucao em BI}%
  \ValidarValorObrigatorio{\@solucaoNotificados}{pessoas notificadas da solucao}%
}

\gdef\@mudancaRitoAutoridade{}
\gdef\@mudancaRitoFundamento{}
\gdef\@mudancaRitoNovasProvas{}
\newcommand{\mudancaRitoAutoridade}[1]{\gdef\@mudancaRitoAutoridade{#1}}
\newcommand{\mudancaRitoFundamento}[1]{\gdef\@mudancaRitoFundamento{#1}}
\newcommand{\mudancaRitoNovasProvas}[1]{\gdef\@mudancaRitoNovasProvas{#1}}
\newcommand{\getMudancaRitoAutoridade}{\@mudancaRitoAutoridade}
\newcommand{\getMudancaRitoFundamento}{\@mudancaRitoFundamento}
\newcommand{\getMudancaRitoNovasProvas}{\@mudancaRitoNovasProvas}
\newcommand{\ValidarMudancaRito}{%
  \ValidarValorObrigatorio{\@mudancaRitoAutoridade}{autoridade que decidiu a mudanca de rito}%
  \ValidarValorObrigatorio{\@mudancaRitoFundamento}{fundamento da mudanca de rito}%
  \ValidarValorObrigatorio{\@mudancaRitoNovasProvas}{novas investigacoes ou provas necessarias}%
  \ifdefstring{\@sindRito}{ordinario}{}{%
    \ClassError{sindicancia}{Anexo AA sem conversao para rito ordinario}{Execute \string\converterRitoParaOrdinario antes de gerar a certidao.}%
  }%
  \ValidarValorObrigatorio{\@sindHistoricoTransicoes}{historico da mudanca de rito}%
}

% ------------------------------------------------------------
% 9. DOCUMENTOS PROCEDIMENTAIS ADICIONAIS
% ------------------------------------------------------------
\gdef\@relatorioComplementarDiligencias{}
\gdef\@relatorioComplementarConclusoes{}
\gdef\@relatorioComplementarResultado{ratifica}
\newcommand{\relatorioComplementarDiligencias}[1]{\gdef\@relatorioComplementarDiligencias{#1}}
\newcommand{\relatorioComplementarConclusoes}[1]{\gdef\@relatorioComplementarConclusoes{#1}}
\newcommand{\relatorioComplementarResultado}[1]{\gdef\@relatorioComplementarResultado{#1}}
\newcommand{\getRelatorioComplementarDiligencias}{\@relatorioComplementarDiligencias}
\newcommand{\getRelatorioComplementarConclusoes}{\@relatorioComplementarConclusoes}
\newcommand{\getRelatorioComplementarResultado}{\@relatorioComplementarResultado}
\newcommand{\ValidarRelatorioComplementar}{%
  \ValidarValorObrigatorio{\@relatorioComplementarDiligencias}{diligencias complementares realizadas}%
  \ValidarValorObrigatorio{\@relatorioComplementarConclusoes}{conclusoes do relatorio complementar}%
}

\gdef\@notificacaoSolucaoNumero{}
\gdef\@notificacaoSolucaoDestinatario{}
\gdef\@notificacaoSolucaoResumo{}
\gdef\@notificacaoSolucaoPrazoRecursoDias{5}
\gdef\@notificacaoSolucaoFormaCiencia{}
\newcommand{\notificacaoSolucaoNumero}[1]{\gdef\@notificacaoSolucaoNumero{#1}}
\newcommand{\notificacaoSolucaoDestinatario}[1]{\gdef\@notificacaoSolucaoDestinatario{#1}}
\newcommand{\notificacaoSolucaoResumo}[1]{\gdef\@notificacaoSolucaoResumo{#1}}
\newcommand{\notificacaoSolucaoPrazoRecursoDias}[1]{\gdef\@notificacaoSolucaoPrazoRecursoDias{#1}}
\newcommand{\notificacaoSolucaoFormaCiencia}[1]{\gdef\@notificacaoSolucaoFormaCiencia{#1}}
\newcommand{\getNotificacaoSolucaoNumero}{\@notificacaoSolucaoNumero}
\newcommand{\getNotificacaoSolucaoDestinatario}{\@notificacaoSolucaoDestinatario}
\newcommand{\getNotificacaoSolucaoResumo}{\@notificacaoSolucaoResumo}
\newcommand{\getNotificacaoSolucaoPrazoRecursoDias}{\@notificacaoSolucaoPrazoRecursoDias}
\newcommand{\getNotificacaoSolucaoFormaCiencia}{\@notificacaoSolucaoFormaCiencia}
\newcommand{\ValidarNotificacaoSolucao}{%
  \ValidarValorObrigatorio{\@notificacaoSolucaoNumero}{numero da notificacao da solucao}%
  \ValidarValorObrigatorio{\@notificacaoSolucaoDestinatario}{destinatario da notificacao da solucao}%
  \ValidarValorObrigatorio{\@notificacaoSolucaoResumo}{resumo da solucao}%
  \ifnum\@notificacaoSolucaoPrazoRecursoDias<5\relax
    \ClassError{sindicancia}{Prazo recursal inferior a cinco dias uteis}{O art. 40 assegura cinco dias uteis.}%
  \fi
}

% ------------------------------------------------------------
% 10. RECURSOS E DEPENDENCIAS DE OUTRAS NORMAS
% ------------------------------------------------------------
\gdef\@sindRecursosRegistrados{}
\gdef\@sindDependenciasNormativas{}
\gdef\@sindPrazosRegistrados{}
\newcommand{\registrarRecurso}[4]{%
  \ValidarValorObrigatorio{#1}{tipo do recurso}%
  \ValidarValorObrigatorio{#2}{data do recurso}%
  \ValidarValorObrigatorio{#3}{recorrente}%
  \ValidarValorObrigatorio{#4}{situacao ou resultado do recurso}%
  \g@addto@macro\@sindRecursosRegistrados{\item \textbf{#1}, em #2, por #3: #4.}%
}
\newcommand{\registrarDependenciaNormativa}[3]{%
  \ValidarValorObrigatorio{#1}{identificacao da norma dependente}%
  \ValidarValorObrigatorio{#2}{tema regulado pela norma dependente}%
  \ValidarValorObrigatorio{#3}{observacao de aplicabilidade normativa}%
  \g@addto@macro\@sindDependenciasNormativas{\item \textbf{#1} -- #2. #3.}%
}
\newcommand{\registrarPrazo}[5]{%
  \ValidarValorObrigatorio{#1}{ato sujeito a prazo}%
  \ValidarValorObrigatorio{#2}{marco inicial do prazo}%
  \ValidarValorObrigatorio{#3}{vencimento do prazo}%
  \ValidarValorObrigatorio{#4}{criterio de contagem do prazo}%
  \ValidarValorObrigatorio{#5}{situacao do prazo}%
  \g@addto@macro\@sindPrazosRegistrados{\item \textbf{#1}: de #2 a #3; #4; situação: #5.}%
}
\newcommand{\getRecursosRegistrados}{\@sindRecursosRegistrados}
\newcommand{\getDependenciasNormativas}{\@sindDependenciasNormativas}
\newcommand{\getPrazosRegistrados}{\@sindPrazosRegistrados}
\newcommand{\QuadroRecursosRegistrados}{%
  \ifcsempty{@sindRecursosRegistrados}{Nenhum recurso registrado.}{\begin{itemize}\@sindRecursosRegistrados\end{itemize}}%
}
\newcommand{\QuadroDependenciasNormativas}{%
  \ifcsempty{@sindDependenciasNormativas}{Nenhuma dependência normativa registrada.}{\begin{itemize}\@sindDependenciasNormativas\end{itemize}}%
}
\newcommand{\QuadroPrazosRegistrados}{%
  \ifcsempty{@sindPrazosRegistrados}{Nenhum prazo registrado.}{\begin{itemize}\@sindPrazosRegistrados\end{itemize}}%
}

% ------------------------------------------------------------
% 11. APLICABILIDADE DOS MODELOS AO ESTADO PROCESSUAL
% ------------------------------------------------------------
\newcommand{\sind@ExigirSindicado}[1]{%
  \ifdefstring{\@sindNatureza}{processual}{}{%
    \ClassError{sindicancia}{Documento #1 exige natureza processual}{Converta a apuracao para processual e registre a certidao correspondente.}%
  }%
  \ifdefstring{\@sindPartePrincipal}{sindicado}{}{%
    \ClassError{sindicancia}{Documento #1 exige sindicado}{Defina a parte principal como sindicado.}%
  }%
}

\newcommand{\ValidarEstadoDocumento}[1]{%
  \ValidarRegimeNormativo
  \ValidarEstadoSindicancia
  \ifstrequal{#1}{notificacao-previa}{\sind@ExigirSindicado{#1}}{}%
  \ifstrequal{#1}{notificacao-novo-sindicado}{\sind@ExigirSindicado{#1}}{}%
  \ifstrequal{#1}{notificacao-diligencias-complementares}{\sind@ExigirSindicado{#1}}{}%
  \ifstrequal{#1}{comparecimento-sindicado}{\sind@ExigirSindicado{#1}}{}%
  \ifstrequal{#1}{termo-inquiricao-sindicado}{\sind@ExigirSindicado{#1}}{}%
  \ifstrequal{#1}{vista-alegacoes-finais}{\sind@ExigirSindicado{#1}}{}%
  \ifstrequal{#1}{termo-consentimento-prontuario}{\sind@ExigirSindicado{#1}}{}%
}

% Estado inicial dos auxiliares dinamicos.
\novaAcareacaoDinamica
